\documentclass[11pt]{article}

\usepackage{amsmath, amssymb}
\usepackage{geometry}
\usepackage{booktabs}
\usepackage{hyperref}
\usepackage{parskip}
\usepackage{titlesec}
\usepackage{array}
\usepackage{xcolor}

\geometry{margin=1.25in}

\title{Neuron Model Options for the Distribution of Oscillators\\
\large C.\ elegans Locomotion Circuit}
\author{Distribution of Oscillators Project}
\date{\today}

\begin{document}

\maketitle

\begin{abstract}
The current piecewise-linear (Yuval) neuron model exhibits intrinsic bistability that creates
network-level locking and prevents reliable oscillations. This document surveys three alternative
models---FitzHugh--Nagumo, Izhikevich, and a graded potential model---evaluating each against
the project's core requirements: few free parameters, no pathological bistability, and biological
appropriateness for \textit{C.~elegans} locomotion.
\end{abstract}

% ─────────────────────────────────────────────────────────────────────────────
\section{Motivation}

The current model uses a piecewise-linear approximation to the membrane voltage dynamics with
three key voltage levels: resting potential $L$, threshold $T$, and plateau $H$.
This produces a neuron with \emph{two stable fixed points}---rest ($V = L$) and plateau
($V = H$)---separated by an unstable threshold at $T$.

In isolation this bistability is useful: it distinguishes oscillator cells (driven above $T$)
from passive cells (below $T$), which is central to the Distribution of Oscillators framework.
However, at the network level it creates two serious problems.

\begin{enumerate}
    \item \textbf{Plateau locking.} Any excitatory connection strong enough to push a target
          neuron above $T$ locks it at the plateau indefinitely. The neuron has no intrinsic
          restoring force back below $T$ once it crosses that boundary.
    \item \textbf{Cascade freezing.} Because excitatory neurons in the six-segment network
          (types 1--3) receive no inhibitory synaptic input, a single locking event can
          propagate through the excitatory subnetwork and freeze the entire circuit.
\end{enumerate}

The fundamental requirement for a replacement model is therefore: a neuron that can oscillate
or be driven to fire, but which \emph{returns to rest naturally} when synaptic drive is
removed, without requiring an external hyperpolarizing current.

% ─────────────────────────────────────────────────────────────────────────────
\section{Model Descriptions}

\subsection{FitzHugh--Nagumo (FHN)}

\begin{align}
    \frac{dv}{dt} &= v - \frac{v^3}{3} - w + I \label{eq:fhn_v} \\
    \frac{dw}{dt} &= \frac{1}{\tau}(v + a - bw) \label{eq:fhn_w}
\end{align}

\textbf{Parameters:} $a$, $b$, $\tau$, $I$ (4 total).

The variable $v$ represents membrane voltage and $w$ is a slower recovery variable analogous
to potassium channel gating. The cubic nullcline of $v$ produces a single stable fixed point
for $I$ below the bifurcation and a limit cycle for $I$ above it. Crucially, the rest state
is \emph{globally attracting} when $I$ is removed---there is no plateau to lock onto.

The oscillator/passive distinction maps cleanly onto $I$ relative to the Hopf bifurcation
point, preserving the conceptual structure of the OscillatorComb framework.

\subsection{Izhikevich}

\begin{align}
    \frac{dv}{dt} &= 0.04v^2 + 5v + 140 - u + I \label{eq:izh_v} \\
    \frac{du}{dt} &= a(bv - u) \label{eq:izh_u}
\end{align}
with the reset rule: if $v \geq 30$\,mV then $v \leftarrow c$, $u \leftarrow u + d$.

\textbf{Parameters:} $a$, $b$, $c$, $d$ (4 total, plus $I$).

The Izhikevich model is a quadratic integrate-and-fire neuron augmented with an adaptation
variable $u$. The four parameters can be tuned to reproduce a wide variety of firing patterns
(tonic spiking, bursting, chattering, adaptation) without changing the model structure.
The reset rule replaces the plateau mechanism entirely, so there is no intrinsic bistability.

\subsection{Graded Potential Model}

\begin{equation}
    \tau \frac{dV}{dt} = -V + \tanh\!\bigl(\beta (I_{\mathrm{syn}} + I_{\mathrm{app}})\bigr)
    \label{eq:graded}
\end{equation}

\textbf{Parameters:} $\tau$, $\beta$ (2 total per cell, plus $I_{\mathrm{app}}$).

This is a rate-coded model in which $V$ represents the mean firing rate or graded membrane
potential rather than an instantaneous voltage. The sigmoidal transfer function $\tanh$ bounds
the output and provides a smooth, globally stable response. There is no threshold crossing
and no bistability by construction.

This model is the most commonly used in published \textit{C.~elegans} locomotion simulations
\cite{karbowski2008, izquierdo2013, olivares2021} because of its biological basis: the
majority of \textit{C.~elegans} neurons are thought to be \textbf{non-spiking}, communicating
via graded potentials rather than action potentials \cite{goodman1998}.

% ─────────────────────────────────────────────────────────────────────────────
\section{Comparison}

\begin{center}
\renewcommand{\arraystretch}{1.4}
\begin{tabular}{p{3.8cm} p{2.2cm} p{2.2cm} p{2.2cm} p{2.2cm}}
\toprule
& \textbf{Yuval PL} & \textbf{FHN} & \textbf{Izhikevich} & \textbf{Graded} \\
\midrule
Free parameters        & 8        & 4         & 4          & 2 \\
Bistable               & Yes      & No        & No         & No \\
Intrinsic oscillations & Via SF   & Yes       & Yes        & No (network only) \\
Plateau locking risk   & High     & None      & None       & None \\
OscillatorComb compatible & Yes  & Yes       & Yes        & Requires redesign \\
Biophysical realism    & Low      & Low       & Medium     & Medium--High \\
C.\ elegans appropriate & Unclear & Unclear   & Unclear    & \textbf{Yes} \\
Computational cost     & Low      & Low       & Low        & Very low \\
Mathematical tractability & High  & High      & Medium     & High \\
\bottomrule
\end{tabular}
\end{center}

% ─────────────────────────────────────────────────────────────────────────────
\section{Tradeoffs}

\subsection{FitzHugh--Nagumo}

FHN is the natural successor to the current model. It retains the two-variable structure
(voltage + recovery, analogous to $V$ + SF), has the same parameter count as the project
requires, and its bifurcation structure is well-understood analytically. The oscillator/passive
distinction is controlled purely by $I$, so the OscillatorComb designation maps directly onto
the applied current without changing the model equations.

The main weakness is that FHN is phenomenological and dimensionless---voltage and time are
not in physical units without rescaling. If the model needs to interface with biophysical
measurements (e.g.\ calcium imaging timescales), a rescaling step must be added.

\subsection{Izhikevich}

The Izhikevich model operates in millivolts and milliseconds, making it easier to set
biophysically plausible synaptic reversal potentials and time constants. The discontinuous
reset rule is the main complication: it requires careful handling in stiff ODE solvers
(event detection) and can produce chattering near the threshold in some parameter regimes,
similar to the piecewise-linear boundary problem in the current model.

The flexibility of the four parameters is both a strength and a weakness---the parameter
space is large and the relationship between parameters and firing pattern is not always
intuitive.

\subsection{Graded Potential}

The graded model is the strongest biological choice. Non-spiking, graded-potential transmission
is well-documented in \textit{C.~elegans} motor circuits \cite{goodman1998}, and most
published models of \textit{C.~elegans} locomotion use this framework. It has the fewest
parameters and is the least computationally expensive.

The key conceptual tradeoff is that oscillations are entirely a \emph{network} property in
this model: no single neuron oscillates in isolation. The OscillatorComb framework, which
designates individual cells as oscillators, would need to be reinterpreted---oscillator cells
might instead be those with higher $\beta$ or $I_{\mathrm{app}}$, or those in particular
network positions, rather than cells with intrinsic limit-cycle dynamics.

% ─────────────────────────────────────────────────────────────────────────────
\section{Recommendation Summary}

\begin{itemize}
    \item \textbf{Preserve the current framework with minimal change:} use
          \textbf{FitzHugh--Nagumo}. Same parameter count, no bistability, OscillatorComb
          compatible.
    \item \textbf{Prioritise biological fidelity to \textit{C.~elegans}:} use the
          \textbf{graded potential model}. Fewest parameters, most appropriate neuroscience,
          but requires reconceptualising what an ``oscillator cell'' means.
    \item \textbf{Avoid Izhikevich} unless the spiking timescale in physical units is
          specifically required, as the reset rule reintroduces a variant of the
          threshold-crossing problem.
\end{itemize}

% ─────────────────────────────────────────────────────────────────────────────
\begin{thebibliography}{9}

\bibitem{karbowski2008}
Karbowski, J.\ \& Bhatt, D. (2008).
Computational model of \textit{C.~elegans} locomotion.
\textit{PLOS Computational Biology}.

\bibitem{izquierdo2013}
Izquierdo, E.~J.\ \& Beer, R.~D. (2013).
Connecting a connectome to behavior: an ensemble of neuroanatomical models of
\textit{C.~elegans} klinotaxis.
\textit{PLOS Computational Biology}, 9(2).

\bibitem{olivares2021}
Olivares, E.~O.\ et al. (2021).
Potential role of a ventral nerve cord central pattern generator in forward and backward
locomotion in \textit{C.~elegans}.
\textit{Network Neuroscience}.

\bibitem{goodman1998}
Goodman, M.~B.\ et al. (1998).
Active currents regulate sensitivity and dynamic range in \textit{C.~elegans} neurons.
\textit{Neuron}, 20(4), 763--772.

\end{thebibliography}

\end{document}
